\chapter{Les histoires que l'on garde}
\label{chap:histoires-que-lon-garde}

Avec Rielle, Kamatsu avait découvert qu'il existait différentes manières de
conserver ce que l'on rencontrait sur les routes.

Il avait continué à dessiner après son départ, pas toujours avec beaucoup de
réussite, mais suffisamment pour remplir peu à peu quelques feuilles de
visages, de paysages et de détails aperçus au cours de leurs voyages. À côté de
certains dessins apparaissaient désormais quelques mots : le nom d'un endroit,
celui d'une personne, une phrase entendue au détour d'une conversation ou
simplement une idée qu'il ne voulait pas oublier.

Puis Thom Merilin rejoignit la caravane.

\scenebreak

Il arriva un après-midi pluvieux, avec peu d'affaires à l'exception d'un
instrument soigneusement protégé des intempéries et d'un sac qui, selon lui,
contenait tout ce dont un homme raisonnable pouvait avoir besoin. Personne ne
sut jamais précisément ce que Thom considérait comme raisonnable.

Le soir même, il s'installa près du feu et sortit son instrument. Les premières
notes attirèrent quelques voyageurs, puis d'autres, jusqu'à former autour de
lui un petit groupe. Il commença par jouer seul avant d'entonner une chanson
dont certains connaissaient manifestement les paroles. Quelques voix se
joignirent à la sienne tandis que les autres se contentaient d'écouter.

Une chanson en suivit une autre, puis Thom finit par poser son instrument pour
raconter une histoire.

Kamatsu ne le quitta pratiquement pas des yeux. Ce n'était pas seulement ce que
racontait le ménestrel qui l'intéressait, mais sa manière de le faire. Thom
changeait légèrement de voix lorsqu'un personnage parlait, ralentissait à
certains moments, accélérait à d'autres et pouvait laisser passer plusieurs
secondes de silence avant une phrase importante.

Et surtout, les gens écoutaient.

Pendant quelques instants, la route, les chariots et le camp semblaient avoir
disparu. Lorsque Thom termina, Kamatsu avait presque oublié qu'il était assis
près d'un feu.

\scenebreak

Les jours suivants, Kamatsu fredonna régulièrement les mélodies qu'il avait
entendues. Certaines lui revenaient facilement, d'autres beaucoup moins. L'une
d'elles lui résistait particulièrement et, depuis le matin, il répétait
obstinément le même passage.

Au bout d'un moment, Kael leva les yeux de ce qu'il était en train de faire.

— Ce n'est pas comme ça.

Kamatsu s'arrêta.

— Si.

— Non.

— Si.

Kael secoua la tête.

— Tu changes la fin.

— C'est toi qui l'as mal retenue.

— Je l'ai entendue aussi.

Kamatsu reprit exactement le même passage, mais Kael l'interrompit presque
immédiatement.

— Voilà. Là.

— Quoi, là ?

— Tu chantes faux.

Kamatsu le regarda.

— C'est toi qui chantes faux.

— Je ne chante pas faux.

— Tu ne chantes même pas.

Kael haussa les épaules.

— Justement.

Une voix intervint derrière eux.

— C'est effectivement une excellente manière d'éviter de chanter faux.

Les deux garçons se retournèrent. Thom les avait entendus.

Le ménestrel regarda Kamatsu.

— Recommence.

Kamatsu hésita avant de reprendre la chanson depuis le début. Thom ne dit rien
et le laissa aller jusqu'au refrain.

— Tu l'as entendue combien de fois ?

Kamatsu réfléchit.

— Deux.

Thom haussa légèrement les sourcils.

— Et tu te souviens de tout ça ?

— J'ai oublié un passage.

Kael intervint immédiatement.

— Deux.

Kamatsu lui lança un regard.

— Un.

— Deux.

Thom sourit.

— On va commencer par celui sur lequel vous êtes d'accord.

Ce fut le début de leurs leçons.

\scenebreak

Thom lui apprit d'abord à mieux utiliser sa voix : respirer au bon moment, ne
pas essayer de chanter plus fort simplement parce qu'il voulait être entendu,
trouver le rythme et écouter les autres lorsqu'il chantait avec eux plutôt que
de chercher à les couvrir.

Il lui montra également quelques rudiments de son instrument. Kamatsu apprit
quelques notes, puis quelques accords simples, avec suffisamment d'erreurs pour
que les premiers essais soient probablement plus instructifs pour lui que
plaisants pour ceux qui se trouvaient à proximité. Kael eut d'ailleurs
rapidement tendance à trouver quelque chose à faire à l'autre bout du camp
lorsque Kamatsu s'entraînait.

Kamatsu prétendit ne pas le remarquer.

Mais Thom ne lui apprenait pas seulement la musique.

Un soir, après lui avoir fait répéter une chanson, il lui demanda de raconter
une histoire.

— Laquelle ?

Thom haussa les épaules.

— N'importe laquelle.

Kamatsu choisit une anecdote arrivée quelques jours auparavant et la raconta.
Lorsqu'il eut terminé, Thom attendit un instant.

— C'est tout ?

— Oui.

— C'était très mauvais.

Kamatsu le regarda, vexé.

Thom sourit.

— Mais l'histoire était bonne.

— Alors pourquoi c'était mauvais ?

— Parce que tu me l'as récitée.

Kamatsu fronça les sourcils.

— Je viens de la raconter.

— Non. Tu m'as expliqué ce qui s'est passé.

Thom se pencha légèrement vers lui.

— Fais-moi y être.

Kamatsu resta silencieux. Cette distinction lui rappelait quelque chose :
Rielle lui avait appris qu'il existait une différence entre regarder quelque
chose et réellement l'observer. Thom venait peut-être de lui montrer qu'il
existait la même différence entre dire ce qui s'était passé et raconter une
histoire.

\scenebreak

Les leçons continuèrent. Kamatsu apprit à choisir les détails qui comptaient et
à laisser de côté ceux qui n'apportaient rien, à ralentir lorsqu'il voulait que
les autres attendent avec lui et, parfois, à se taire plutôt que de remplir
chaque silence.

Surtout, Thom lui apprit à regarder ceux qui l'écoutaient.

Un soir, il demanda à Kamatsu pourquoi il observait autant son public lorsqu'il
chantait.

— Pour voir s'ils aiment bien.

— Mauvaise réponse.

Kamatsu attendit.

— Regarde-les pour savoir s'ils écoutent.

Kamatsu observa les gens autour du feu. Certains parlaient entre eux, d'autres
mangeaient, une femme réparait un vêtement, Kael s'occupait de son arc et
Bartiméus dormait près des flammes sans manifester le moindre intérêt
artistique.

— Ils ne m'écoutent pas.

— Alors donne-leur une raison de le faire.

Kamatsu regarda de nouveau les voyageurs avant de recommencer. Cette fois, il
chanta moins fort, mais différemment, en prenant davantage son temps.

Lorsqu'arriva le refrain, deux personnes avaient cessé leur conversation. Une
troisième tourna bientôt la tête vers lui et, un peu plus tard, Kael cessa
finalement de s'occuper de son arc.

Ce soir-là, personne n'applaudit particulièrement lorsqu'il termina, mais Thom
souriait.

— Tu as vu ?

Kamatsu acquiesça.

Ils avaient écouté.

\scenebreak

Quelques jours plus tard, Kamatsu osa chanter devant toute la caravane. Cette
fois, il ne s'installa pas discrètement à côté de Thom pour s'entraîner, mais
prit sa place près du feu.

Les premières paroles furent hésitantes. Sa voix trembla légèrement et il
manqua presque son entrée sur le deuxième couplet, puis il aperçut les visages
autour de lui et se souvint de la leçon.

Il cessa de penser à sa voix.

Il raconta la chanson.

Progressivement, les conversations cessèrent. Quelqu'un commença à battre la
mesure, puis une voix reprit le refrain, bientôt suivie d'une autre. Lorsque
Kamatsu arriva au dernier, plusieurs voyageurs chantaient avec lui.

La dernière note disparut dans les bruits du camp et Kamatsu chercha Thom du
regard. Le ménestrel acquiesça simplement.

— Maintenant, ils écoutaient.

Kamatsu sourit.

Kael, assis un peu plus loin, attendit quelques secondes.

— Tu as quand même changé un mot.

Kamatsu tourna lentement la tête vers lui.

— Non.

— Si.

— Personne ne l'a remarqué.

— Moi, si.

Kamatsu attrapa ce qu'il avait sous la main et le lui lança. Kael l'esquiva en
riant tandis que Bartiméus ouvrait un œil pour observer les deux garçons avant
de le refermer.

L'incident ne méritait manifestement pas son intervention.

\scenebreak

Le jour où Thom annonça qu'il allait quitter la caravane, Kamatsu vint le
retrouver pendant qu'il préparait ses affaires. Il avait passé suffisamment de
soirées à ses côtés pour connaître plusieurs de ses chansons, mais pas assez à
son goût.

— Je ne connais toujours pas toutes vos chansons.

— Heureusement.

Kamatsu le regarda, surpris.

— Pourquoi ?

Thom referma son sac avant de se tourner vers lui.

— Parce que si tu les connaissais toutes, tu passerais peut-être ton temps à
chanter les miennes.

Il désigna les quelques feuilles que Kamatsu transportait désormais
régulièrement avec lui.

— Écoute les histoires des autres. Apprends leurs chansons. Note ce que tu ne
veux pas oublier.

Puis il sourit.

— Mais garde toujours une place pour les tiennes.

Kamatsu baissa les yeux vers ses feuilles.

— Je n'ai pas d'histoires.

Thom eut un petit rire.

— Tu voyages sur les routes depuis des années et tu crois ne pas avoir
d'histoires ?

Kamatsu réfléchit tandis que Thom passait son sac sur son épaule.

— Si tu passes ta vie à raconter les histoires des autres, qui racontera les
tiennes ?

Cette fois, Kamatsu ne répondit pas. Thom n'attendait probablement pas de
réponse.

Il aperçut alors Kael quelques pas plus loin.

— Et toi, continue de lui dire quand il change les paroles.

Kael acquiesça avec beaucoup trop de sérieux.

— Je peux faire ça.

— Kael...

Thom partit en riant.

\scenebreak

Le soir suivant, quelqu'un demanda à Kamatsu de chanter. Après une brève
hésitation, il s'installa près du feu et choisit l'une des chansons que Thom
lui avait enseignées.

Les paroles étaient les mêmes, à peu près, et la mélodie aussi. Mais lorsqu'il
commença à chanter, Kamatsu ne chercha pas à reproduire la voix de Thom, ses
gestes ou sa manière de raconter.

Il regarda ceux qui l'écoutaient.

Puis il raconta l'histoire à sa manière.

Pour la première fois, ce n'était plus une chanson de Thom que Kamatsu essayait
de reproduire.

C'était une chanson que Kamatsu racontait.